\documentclass[11pt]{article}
\usepackage[margin=1in]{geometry}
\usepackage{booktabs}
\usepackage{xcolor}
\usepackage{enumitem}
\usepackage{amsmath}
\usepackage{hyperref}
\hypersetup{colorlinks=true, linkcolor=blue, urlcolor=blue}

\title{DendriNet Kernel Optimization: Work Log and Results}
\author{Ryan Whalen}
\date{July 2026}

\newcommand{\what}[1]{\textbf{What it is:} #1}
\newcommand{\impact}[1]{\textbf{Impact:} #1}

\begin{document}
\maketitle

\section*{Context}
This document logs a profile-driven optimization pass on the DendriNet codebase, targeting
the sparse ``indexed'' synapse used by the excitation--inhibition population network
(config \texttt{sglimit5\_h100\_cifar10\_ei\_1layer\_bf33\_10000e1250i}). All work is at
the \emph{Python / PyTorch} level---no CUDA or Triton kernels were written. The guiding
discipline: every change must be \emph{explicitly visible}, \emph{reversible}, and a
\emph{strict improvement with no behavior change by default} (safe to merge into
\texttt{develop}).

\begin{center}
\fbox{\parbox{0.92\textwidth}{\centering
\textbf{Headline result.} A single custom-autograd kernel made the sparse-synapse
\emph{backward} pass \textbf{7.6$\times$ faster} and \textbf{28$\times$ lighter on memory},
cutting whole-step wall-clock \textbf{$\sim$5$\times$} (2796\,ms $\to$ 554\,ms), with
\textbf{bit-identical training} under mixed precision. Gated behind a default-off flag.
}}
\end{center}

\section{Investigation \& Diagnosis}

\begin{enumerate}[leftmargin=*]

\item \textbf{Traced the mystery profiling report.}\\
\what{Figured out which model produced an \texttt{.nsys-rep} file whose name we were handed.}\\
\impact{Found it came from a separate profiling harness (``ProfilePilot,'' another user's
setup), and that the config---despite the name saying ``1layer''---is actually a
\emph{5-layer feedforward} EI network on CIFAR-10. Prevented us from optimizing the wrong thing.}

\item \textbf{Mapped what the model actually computes.}\\
\what{Traced every PyTorch operation the model runs in its forward/backward pass.}\\
\impact{Confirmed there are \emph{no} CUDA/Triton kernels and \emph{no} attention---just
gather, multiply, sum, matmul, and elementwise ops. Established the ``menu'' of things that
could be optimized.}

\item \textbf{Built a shared-primitive leverage table.}\\
\what{Ranked operations by how many models use them (breadth), so one fix could help many configs.}\\
\impact{Identified the high-leverage shared ops (weight transform, sparse gather, block-linear
aggregation) and the principle: target code that is \emph{both} widely-shared \emph{and}
inefficient.}

\item \textbf{Profiled first---and killed our initial hypothesis.}\\
\what{Put a stopwatch on every GPU operation to see where time \emph{actually} goes.}\\
\impact{Our first idea (caching the weight transform) turned out to be \textbf{$<$0.3\%} of
runtime---invisible. Profiling redirected us to the real cost before we wasted effort. Key
lesson: the slow part is usually not the code you can see.}

\item \textbf{Located the true bottleneck.}\\
\what{The 84\% hotspot was the \emph{backward} (gradient) pass of a sparse gather---a step
PyTorch auto-generates, not code anyone wrote.}\\
\impact{Intuition: thousands of outputs read the same input, so during learning every
correction has to be ``added back'' onto shared inputs. PyTorch's default did this by making
them queue one-at-a-time---running at only 3\% of the GPU's memory-bandwidth ceiling.}

\end{enumerate}

\section{The Fix}

\begin{enumerate}[leftmargin=*, resume]

\item \textbf{Isolated microbenchmark with proper metrics.}\\
\what{A standalone script timing just the gather op, measuring achieved memory bandwidth vs.\
the GPU's peak (the ``roofline'' test).}\\
\impact{Proved the op was memory-bound and stalling---quantified headroom without the noise of
a full model run.}

\item \textbf{Prototyped the optimized kernel.}\\
\what{A custom \texttt{autograd.Function} that overrides the slow default backward with an
explicit grouped scatter (\texttt{index\_add\_}), and recomputes a large temporary instead of
storing it.}\\
\impact{\textbf{7.6$\times$ faster backward} (303.7\,ms $\to$ 39.8\,ms) and \textbf{28$\times$
less memory} (5.64\,GB $\to$ 0.20\,GB) at production scale.}

\item \textbf{Verified correctness rigorously.}\\
\what{Compared the new kernel against the original in single (fp32) and double (fp64) precision.}\\
\impact{Outputs and weight-gradients \emph{bit-identical}; the input-gradient differs only by
floating-point rounding (exact in fp64), far below the precision training uses. Confirmed
\emph{not a bug}.}

\item \textbf{Shipped a production, memory-safe module.}\\
\what{Turned the prototype into a real module (\texttt{fused\_gather.py}), chunked internally so
it never allocates the huge intermediate tensor.}\\
\impact{Runs at the real \texttt{out=90000} scale that previously caused out-of-memory errors.}

\item \textbf{Integrated behind a default-off flag.}\\
\what{Added \texttt{use\_fused\_gather} (default \texttt{False}) to the synapse; when off, the
original code path runs unchanged.}\\
\impact{Merge-safe: nothing changes for existing runs unless explicitly enabled. One-line,
greppable, reversible.}

\item \textbf{Proved whole-model training-equivalence.}\\
\what{Trained the real model twice from an identical start---original vs.\ fused kernel---and
compared the loss trajectory.}\\
\impact{\textbf{Bit-identical} loss every step (\texttt{max diff = 0.0}); final weights matched
to $\sim$$10^{-6}$. The strongest possible correctness guarantee.}

\item \textbf{Plumbed the flag through to YAML.}\\
\what{Wired \texttt{indexed\_use\_fused\_gather: true} from config files down to the synapse
(6 small edits, all defaulting off).}\\
\impact{The optimization can be toggled per-experiment from a config, with zero effect when unset.}

\end{enumerate}

\section{Re-measurement \& Next-Candidate Search}

\begin{enumerate}[leftmargin=*, resume]

\item \textbf{Re-profiled with the fix enabled.}\\
\what{Measured the model \emph{after} the fix to find the new bottleneck (the ranking shifts once
the biggest cost is removed).}\\
\impact{Confirmed \textbf{$\sim$5$\times$ whole-step speedup} (2796\,ms $\to$ 554\,ms). The
remaining work is now \emph{diffuse} (many $\sim$20\% ops) and mildly \emph{launch-bound}
(CPU dispatch $>$ GPU compute).}

\item \textbf{A/B tested \texttt{torch.compile}.}\\
\what{A whole-graph tool that fuses many small ops into fewer, larger kernels.}\\
\impact{\textbf{1.55$\times$ faster} (555\,ms $\to$ 355\,ms) with numerically-faithful gradients
---\emph{but} it uses \emph{more} memory (re-materializing the tensor our kernel removed, causing
OOM on the small GPU) and is not bit-identical. \textbf{Conclusion:} an opt-in speedup for
memory-rich hardware, \emph{not} a strict default. Correctly \emph{rejected} as a default by
measurement.}

\end{enumerate}

\section{Deliverables}

\begin{table}[h]
\centering
\small
\begin{tabular}{@{}lll@{}}
\toprule
\textbf{File} & \textbf{Type} & \textbf{Purpose} \\
\midrule
\texttt{synapse/fused\_gather.py} & new source & the optimized kernel (custom autograd Function) \\
\texttt{synapse/indexed\_sparse.py} & modified & flag + forward branch (only real logic change) \\
\texttt{recurrent/ei\_config.py} & modified & config schema field \\
\texttt{dendritic/synapse\_config.py} & modified & YAML$\leftrightarrow$kwargs mapping \\
\texttt{dendritic/branch\_layer.py} & modified & reads flag onto branch layer \\
\texttt{dendritic/branch\_synapses.py} & modified & passes flag to constructor \\
\midrule
\texttt{profile\_weight\_transform.py} & tool & full-model profiler (find hotspots) \\
\texttt{bench\_kernel.py} & tool & isolated baseline microbenchmark \\
\texttt{bench\_kernel-rzw-optimized.py} & tool & optimized-kernel microbench + correctness \\
\texttt{verify\_fused\_gather.py} & tool & full-model training-equivalence gate \\
\texttt{profile\_postfix.py} & tool & re-profile after fix (find next candidate) \\
\texttt{compile\_ab.py} & tool & torch.compile A/B + faithfulness check \\
\bottomrule
\end{tabular}
\caption{Six source edits (24 insertions, 7 deletions; all default-off) plus six dev tools.}
\end{table}

\section{Results Summary}

\begin{table}[h]
\centering
\small
\begin{tabular}{@{}lrrr@{}}
\toprule
\textbf{Metric} & \textbf{Before} & \textbf{After} & \textbf{Change} \\
\midrule
Sparse backward latency & 303.7\,ms & 39.8\,ms & \textbf{7.6$\times$ faster} \\
Backward bandwidth utilization & 3\% of peak & 20\% of peak & 7$\times$ better \\
Peak memory (the op) & 5.64\,GB & 0.20\,GB & \textbf{28$\times$ less} \\
Whole-step wall-clock & 2796\,ms & 554\,ms & \textbf{$\sim$5$\times$ faster} \\
Training correctness & --- & bit-identical (AMP) & no behavior change \\
\midrule
\multicolumn{4}{@{}l}{\emph{torch.compile (evaluated, not merged as default):} 1.55$\times$ faster but $+$memory, not bit-exact} \\
\bottomrule
\end{tabular}
\caption{Measured on NVIDIA A100 MIG 3g.20gb, dominant shape (out=90000, K=40, batch=128).}
\end{table}

\section*{Method (reusable playbook)}
\begin{enumerate}[leftmargin=*, itemsep=1pt]
\item Pick a candidate op that is shared across many configs (breadth---checked statically first).
\item Profile a representative config to confirm it is actually a large fraction of time.
\item Isolate it in a microbenchmark; write a faster equivalent; prove identical output + gradients.
\item Gate behind a default-off flag; verify whole-model training-equivalence.
\item Re-profile---the ranking reshuffles---and repeat on the new top candidate.
\end{enumerate}

\noindent \emph{Status: sparse-gather optimization complete and merge-ready.
\texttt{torch.compile} logged as an opt-in for memory-rich hardware. Next: Kempner pulse
ingestion.}

\end{document}
